\subsection{\texorpdfstring{$p$}{p}-进 Prony--Hankel 精度塔的格子扩张、条件数账本与 simultaneous de-resonance}\label{subsec:conclusion-padic-prony-lattice-deresonance-precisiontower}
本小节把非阿基米德 Hankel 理论中的四个此前分散的接口压成同一条精度塔主线：
主导链将主子式赋值写成成对碰撞能量，
伴随矩阵给出 sharp 的解坐标估值轮廓，
偏移 Hankel 塔把这一精度税升级为格子扩张指数，
而有限剩余域上的单一单位扰动又能以正概率同时 rigidify 整条主子式塔。

\begin{theorem}[主导链下的 Hankel 精度税等于成对 \texorpdfstring{$p$}{p}-进能量]\label{thm:conclusion-padic-hankel-pairwise-collision-energy-precision-tax}
在定理 \ref{thm:conclusion-nonarch-hankel-tropical-chain-increment} 的设定下，
对每个 \(0\le r\le \kappa-1\) 记
\[
s_r:=v(\det H_r).
\]
则有精确公式
\[
s_r
=
2\sum_{1\le i<j\le r+1}v(a_j-a_i).
\]
特别地，对任意
\[
b\in \mathcal O_K^{r+1}
\]
与线性系统
\[
H_r x=b,
\]
其所有解都满足
\[
x\in \pi^{-s_r}\mathcal O_K^{r+1}.
\]
\end{theorem}
\begin{proof}
由定理 \ref{thm:conclusion-nonarch-hankel-tropical-chain-increment}，
\[
v(\det H_r)
=
2\sum_{j=2}^{r+1}v\!\left(\prod_{i=1}^{j-1}(a_j-a_i)\right).
\]
把乘积赋值展开，即得
\[
s_r
=
2\sum_{j=2}^{r+1}\sum_{i=1}^{j-1}v(a_j-a_i)
=
2\sum_{1\le i<j\le r+1}v(a_j-a_i).
\]
另一方面，对 \(H_r\in M_{r+1}(\mathcal O_K)\) 有
\[
H_r^{-1}=\frac{\operatorname{adj}(H_r)}{\det(H_r)}.
\]
由于 \(\operatorname{adj}(H_r)\) 的条目都在 \(\mathcal O_K\) 中，故任意 \(b\in\mathcal O_K^{r+1}\) 的解都满足
\[
x\in \pi^{-v(\det H_r)}\mathcal O_K^{r+1}
=
\pi^{-s_r}\mathcal O_K^{r+1}.
\]
证毕。
\end{proof}

\begin{theorem}[伴随矩阵的列估值轮廓与可达最优解]\label{thm:conclusion-padic-hankel-adjugate-column-valuation-profile}
设
\[
H\in M_d(\mathcal O_K)
\]
可逆，记
\[
s:=v(\det H),\qquad
A=(A_{ij}):=\operatorname{adj}(H),
\qquad
c_j:=\min_{1\le i\le d}v(A_{ij}).
\]
则对任意右端项
\[
b=(b_1,\dots,b_d)^{\mathsf T}\in K^d,
\qquad
x=H^{-1}b,
\]
都有 sharp 下界
\[
\min_{1\le i\le d}v(x_i)\ \ge\ \min_{1\le j\le d}\bigl(v(b_j)+c_j\bigr)-s.
\]
并且该下界对每一列都是可达的：取标准基向量 \(e_j\) 为右端项，则对应解满足
\[
\min_{1\le i\le d}v\bigl((H^{-1}e_j)_i\bigr)=c_j-s.
\]
\end{theorem}
\begin{proof}
由伴随公式
\[
H^{-1}=\frac{\operatorname{adj}(H)}{\det H}
\]
可得对每个 \(i\),
\[
x_i=\frac{1}{\det H}\sum_{j=1}^{d}A_{ij}b_j.
\]
由非阿基米德不等式，
\[
v(x_i)\ge \min_{1\le j\le d}\bigl(v(A_{ij})+v(b_j)\bigr)-s.
\]
由于 \(v(A_{ij})\ge c_j\) 对一切 \(i,j\) 成立，遂有
\[
v(x_i)\ge \min_{1\le j\le d}\bigl(v(b_j)+c_j\bigr)-s.
\]
再对 \(i\) 取最小值得
\[
\min_i v(x_i)\ge \min_j\bigl(v(b_j)+c_j\bigr)-s.
\]

若取 \(b=e_j\)，则
\[
H^{-1}e_j
\]
恰为 \(H^{-1}\) 的第 \(j\) 列，并且
\[
(H^{-1}e_j)_i=\frac{A_{ij}}{\det H}.
\]
故
\[
\min_i v\bigl((H^{-1}e_j)_i\bigr)
=
\min_i\bigl(v(A_{ij})-s\bigr)
=
c_j-s.
\]
从而下界逐列可达。证毕。
\end{proof}

\begin{theorem}[非阿基米德条件数的精确账本公式]\label{thm:conclusion-padic-hankel-exact-condition-ledger}
在定理 \ref{thm:conclusion-padic-hankel-adjugate-column-valuation-profile} 的设定下，
设剩余域基数为
\[
q:=\bigl|\mathcal O_K/\pi\mathcal O_K\bigr|<\infty,
\]
并取归一化绝对值
\[
|\pi|_\pi=q^{-1}.
\]
在 \(K^d\) 上赋予 sup-范数
\[
|x|_\infty:=\max_i |x_i|_\pi.
\]
则逆矩阵的算子范数满足精确闭式
\[
|H^{-1}|_{\infty\to\infty}
=
q^{\,s-c_{\min}},
\qquad
c_{\min}:=\min_{1\le i,j\le d}v(A_{ij}).
\]
因此，\(H\) 的全部 \(p\)-进不稳定性被单一标量
\[
\kappa_\pi(H):=q^{\,s-c_{\min}}
\]
完全承载；这不是上界，而是精确等号。
\end{theorem}
\begin{proof}
对任意矩阵 \(M=(M_{ij})\in M_d(K)\)，非阿基米德 sup-范数下有
\[
|M|_{\infty\to\infty}=\max_{1\le i,j\le d}|M_{ij}|_\pi.
\]
事实上，对任意 \(x\in K^d\),
\[
|(Mx)_i|_\pi
\le
\max_j |M_{ij}|_\pi\,|x|_\infty,
\]
故
\[
|M|_{\infty\to\infty}\le \max_{i,j}|M_{ij}|_\pi.
\]
而对达到右端极大的某个条目 \(M_{i_0j_0}\)，取标准基向量 \(e_{j_0}\) 即有
\[
|Me_{j_0}|_\infty\ge |M_{i_0j_0}|_\pi,
\]
故等号成立。

应用于
\[
(H^{-1})_{ij}=\frac{A_{ij}}{\det H},
\]
得到
\[
|H^{-1}|_{\infty\to\infty}
=
\max_{i,j} q^{-v(A_{ij})+s}
=
q^{\,s-\min_{i,j}v(A_{ij})}
=
q^{\,s-c_{\min}}.
\]
证毕。
\end{proof}

\begin{theorem}[偏移 Hankel 塔的格子扩张律]\label{thm:conclusion-shifted-hankel-lattice-expansion-law}
沿用定理 \ref{thm:conclusion-hankel-shift-padic-precision-slope} 的设定，
并写偏移传播矩阵为
\[
\mathcal A\in M_d(\mathcal O_K),\qquad
H_{k_0+r}=H_{k_0}\mathcal A^r\qquad(r\ge 0),
\]
其中
\[
\det(H_{k_0})\neq 0,\qquad \det(\mathcal A)\neq 0.
\]
定义解格
\[
\Lambda_r:=H_{k_0+r}^{-1}\mathcal O_K^d\subset K^d.
\]
若剩余域基数为 \(q<\infty\)，并记
\[
\lambda:=v(\det \mathcal A)\in\mathbf Z_{\ge 0},
\]
则
\[
\Lambda_r\subseteq \Lambda_{r+1}\qquad(r\ge 0),
\]
并且每一步扩张指数精确为
\[
[\Lambda_{r+1}:\Lambda_r]=q^\lambda.
\]
特别地，
\[
\Lambda_r=\Lambda_0\qquad(\forall r\ge 0)
\iff
\lambda=0.
\]
\leanverified{paper\_conclusion\_shifted\_hankel\_lattice\_expansion\_seeds}
\leanverified{paper\_conclusion\_shifted\_hankel\_lattice\_expansion\_package}
\end{theorem}
\begin{proof}
由
\[
H_{k_0+r+1}=H_{k_0+r}\mathcal A
\]
可得
\[
H_{k_0+r+1}^{-1}=\mathcal A^{-1}H_{k_0+r}^{-1},
\]
从而
\[
\Lambda_{r+1}=\mathcal A^{-1}\Lambda_r.
\]
由于 \(\mathcal A\in M_d(\mathcal O_K)\)，有
\[
\mathcal A\mathcal O_K^d\subseteq \mathcal O_K^d,
\]
故
\[
\mathcal O_K^d\subseteq \mathcal A^{-1}\mathcal O_K^d,
\]
进而推出
\[
\Lambda_r\subseteq \Lambda_{r+1}.
\]

再以 \(H_{k_0+r}\) 的乘法同构把商群 \(\Lambda_{r+1}/\Lambda_r\) 识别为
\[
\mathcal A^{-1}\mathcal O_K^d/\mathcal O_K^d.
\]
对 \(\mathcal A\) 取 Smith 标准形
\[
U\mathcal A V=\operatorname{diag}(\pi^{\lambda_1},\dots,\pi^{\lambda_d}),
\qquad \lambda_i\ge 0,\qquad \sum_i \lambda_i=\lambda.
\]
于是
\[
\mathcal A^{-1}\mathcal O_K^d/\mathcal O_K^d
\cong
\bigoplus_{i=1}^{d}\pi^{-\lambda_i}\mathcal O_K/\mathcal O_K,
\]
其基数为
\[
\prod_{i=1}^{d} q^{\lambda_i}=q^\lambda.
\]
故
\[
[\Lambda_{r+1}:\Lambda_r]=q^\lambda.
\]
最后，
\[
[\Lambda_{r+1}:\Lambda_r]=1
\iff
q^\lambda=1
\iff
\lambda=0,
\]
而这又等价于 \(\Lambda_{r+1}=\Lambda_r\) 对所有 \(r\) 成立。证毕。
\end{proof}

\leanverified{paper\_conclusion\_shifted\_hankel\_uniform\_integrality\_fourfold\_equivalence}
\begin{theorem}[统一积分性的四重等价]\label{thm:conclusion-shifted-hankel-uniform-integrality-fourfold-equivalence}
在定理 \ref{thm:conclusion-shifted-hankel-lattice-expansion-law} 的设定下，以下四个命题彼此等价：
\begin{enumerate}
  \item 对一切 \(r\ge 0\) 与一切 \(b\in \mathcal O_K^d\)，线性系统
  \[
  H_{k_0+r}x=b
  \]
  都有解
  \[
  x\in \mathcal O_K^d;
  \]
  \item 对一切 \(r\ge 0\)，有
  \[
  \Lambda_r=\mathcal O_K^d;
  \]
  \item 有
  \[
  v(\det H_{k_0})=0,\qquad v(\det \mathcal A)=0;
  \]
  \item 有
  \[
  \pi\nmid \det(H_{k_0}),\det(H_{k_0+1}).
  \]
\end{enumerate}
\end{theorem}
\begin{proof}
\((1)\Leftrightarrow(2)\) 由 \(\Lambda_r=H_{k_0+r}^{-1}\mathcal O_K^d\) 的定义立即得到。

\((3)\Leftrightarrow(4)\) 由
\[
\det(H_{k_0+1})=\det(H_{k_0})\det(\mathcal A)
\]
可得
\[
v(\det H_{k_0+1})=v(\det H_{k_0})+v(\det \mathcal A).
\]
故 \(\det(H_{k_0})\) 与 \(\det(H_{k_0+1})\) 同时为单位，当且仅当
\[
v(\det H_{k_0})=0,\qquad v(\det \mathcal A)=0.
\]

现证 \((3)\Rightarrow(2)\)。由 \(v(\det H_{k_0})=0\) 且 \(H_{k_0}\in M_d(\mathcal O_K)\)，其伴随矩阵条目仍在 \(\mathcal O_K\) 中，故
\[
H_{k_0}^{-1}\in M_d(\mathcal O_K),
\]
从而
\[
\Lambda_0=\mathcal O_K^d.
\]
又由 \(v(\det \mathcal A)=0\) 与 \(\mathcal A\in M_d(\mathcal O_K)\)，知 \(\mathcal A\in GL_d(\mathcal O_K)\)，于是
\[
\mathcal A^{-1}\mathcal O_K^d=\mathcal O_K^d.
\]
再由定理 \ref{thm:conclusion-shifted-hankel-lattice-expansion-law} 的公式
\[
\Lambda_{r+1}=\mathcal A^{-1}\Lambda_r
\]
归纳得
\[
\Lambda_r=\mathcal O_K^d\qquad(\forall r\ge 0).
\]

最后证 \((2)\Rightarrow(3)\)。若 \(\Lambda_0=\mathcal O_K^d\)，则 \(H_{k_0}^{-1}\) 为整矩阵，故 \(\det(H_{k_0})\) 为单位，即
\[
v(\det H_{k_0})=0.
\]
若再对一切 \(r\) 有 \(\Lambda_r=\mathcal O_K^d\)，则由定理 \ref{thm:conclusion-shifted-hankel-lattice-expansion-law},
\[
[\Lambda_{r+1}:\Lambda_r]=q^{v(\det \mathcal A)}=1,
\]
故 \(v(\det \mathcal A)=0\)。证毕。
\end{proof}

\begin{corollary}[偏移精度债的二次累积律]\label{cor:conclusion-shifted-hankel-cumulative-precision-debt}
沿用定理 \ref{thm:conclusion-shifted-hankel-lattice-expansion-law} 的记号，记
\[
s_r:=v(\det H_{k_0+r})=s_0+r\lambda,
\qquad
s_0:=v(\det H_{k_0}),
\qquad
\lambda:=v(\det \mathcal A).
\]
定义长度 \(R\) 的累计精度债
\[
\mathcal D_R:=\sum_{r=0}^{R}s_r.
\]
则有精确闭式
\[
\mathcal D_R=(R+1)s_0+\frac{R(R+1)}{2}\lambda.
\]
特别地，
\begin{enumerate}
  \item 若 \(\lambda=0\)，则累计精度债仅线性增长；
  \item 若 \(\lambda>0\)，则累计精度债严格二次增长。
\end{enumerate}
\leanverified{paper\_conclusion\_shifted\_hankel\_precision\_debt\_seeds}
\end{corollary}
\begin{proof}
由
\[
H_{k_0+r}=H_{k_0}\mathcal A^r
\]
取行列式得
\[
\det(H_{k_0+r})=\det(H_{k_0})\det(\mathcal A)^r,
\]
故
\[
s_r=s_0+r\lambda.
\]
逐项求和即得
\[
\mathcal D_R
=
\sum_{r=0}^{R}(s_0+r\lambda)
=
(R+1)s_0+\lambda\sum_{r=0}^{R}r
=
(R+1)s_0+\frac{R(R+1)}{2}\lambda.
\]
最后两条只是对该闭式的直接解释。证毕。
\end{proof}

\begin{theorem}[单一单位扰动对整条 Hankel 主子式塔的 simultaneous de-resonance]\label{thm:conclusion-padic-hankel-simultaneous-deresonance}
沿用命题 \ref{prop:conclusion-nonarch-hankel-random-unit-perturbation} 的设定，
设剩余域
\[
k=\mathcal O_K/(\pi)
\]
有限，\(\card(k)=q\)。
固定某个指标 \(j_0\in\{1,\dots,\kappa\}\)，并作单位扰动
\[
w_{j_0}\longmapsto w_{j_0}(u):=\widehat u\,w_{j_0},
\qquad u\in k^\times,
\]
其中 \(\widehat u\in\mathcal O_K^\times\) 为任一提升。
对每个 \(r=0,\dots,\kappa-1\)，定义
\[
\mathfrak A_r:=
\sum_{\substack{I\in\mathcal S_r\\ j_0\in I}}
\overline{\left(\frac{t_I}{\pi^{\mu_r}}\right)}
\ \in k.
\]
若对一切 \(r\) 都有
\[
\mathfrak A_r\neq 0,
\]
则
\[
\mathbb P_u\Bigl(v(\det H_r(u))=\mu_r\ \text{对所有 }0\le r\le \kappa-1\Bigr)\ge 1-\frac{\kappa}{q}.
\]
\end{theorem}
\begin{proof}
对每个固定的 \(r\)，命题 \ref{prop:conclusion-nonarch-hankel-random-unit-perturbation} 告诉我们：
若 \(\mathfrak A_r\neq 0\)，则坏事件
\[
E_r:=\{u\in k^\times:S_r(u)\equiv 0\pmod{\pi}\}
\]
满足
\[
\mathbb P_u(E_r)\le \frac{1}{q},
\]
并且在补事件 \(E_r^c\) 上有
\[
v(\det H_r(u))=\mu_r.
\]
因此由并集估计，
\[
\mathbb P_u\Bigl(\bigcup_{r=0}^{\kappa-1}E_r\Bigr)
\le
\sum_{r=0}^{\kappa-1}\mathbb P_u(E_r)
\le
\frac{\kappa}{q}.
\]
取补事件即得
\[
\mathbb P_u\Bigl(v(\det H_r(u))=\mu_r\ \forall\,0\le r\le \kappa-1\Bigr)\ge 1-\frac{\kappa}{q}.
\]
证毕。
\end{proof}

\begin{corollary}[有限剩余域上存在一个 global rigidifier]\label{cor:conclusion-padic-hankel-global-rigidifier}
\leanverified{paper\_conclusion\_padic\_hankel\_global\_rigidifier}
在定理 \ref{thm:conclusion-padic-hankel-simultaneous-deresonance} 的设定下，若
\[
q>\kappa,
\]
则存在某个单位
\[
u_\ast\in k^\times
\]
使得整条 Hankel 主子式塔同时达到热带极小值：
\[
v(\det H_r(u_\ast))=\mu_r
\qquad(0\le r\le \kappa-1).
\]
\end{corollary}
\begin{proof}
由定理 \ref{thm:conclusion-padic-hankel-simultaneous-deresonance},
\[
\mathbb P_u\Bigl(v(\det H_r(u))=\mu_r\ \forall\,0\le r\le \kappa-1\Bigr)\ge 1-\frac{\kappa}{q}.
\]
当 \(q>\kappa\) 时，右端严格为正，故该事件非空。于是存在至少一个
\[
u_\ast\in k^\times
\]
满足所述 simultaneous rigidification 条件。证毕。
\end{proof}

综上，\(p\)-进 Hankel 理论在结论层面不再只是若干零散的赋值下界，而形成了一条完整的精度塔链：
主导链把主子式赋值写成成对碰撞能量，
伴随矩阵把这一能量转写成 sharp 的解坐标估值轮廓与精确条件数，
偏移 Hankel 塔则把单层精度税升级为指数为 \(q^{v(\det\mathcal A)}\) 的格子扩张，
而有限剩余域上的单一单位扰动又能以正概率同时 rigidify 全部主子式。
于是，热带主导、Prony 精度税、格子算术与随机去共振，已在同一条非阿基米德精度塔中闭合。

\endinput
